\documentclass[11pt]{article}
\usepackage[a4paper,margin=25mm]{geometry}
\usepackage[T1]{fontenc}
\usepackage{lmodern,amsmath,amssymb,amsthm,booktabs,array}
\usepackage[colorlinks=true,allcolors=blue,breaklinks=true]{hyperref}
\usepackage{microtype}
\hypersetup{pdftitle={Resolving Pozsgay--Wanless's factor-perfection question and qubit decomposition conjecture},
pdfauthor={Haobo Ma and Wenlin Zhang},pdfsubject={Classical and quantum perfect tensor factorization; CC BY 4.0}}
\newtheorem{theorem}{Theorem}
\newtheorem{lemma}[theorem]{Lemma}
\newtheorem{corollary}[theorem]{Corollary}
\newcommand{\OA}{\operatorname{OA}}
\newcommand{\ind}{\mathbf{1}}
\title{Resolving Pozsgay--Wanless's factor-perfection question
and qubit decomposition conjecture\\[4pt]
\large \S5, after equation (5.13), and \S8}
\author{Haobo Ma\thanks{ChronoAI Pte Ltd, Singapore.}\and
Wenlin Zhang\thanks{National University of Singapore; ChronoAI Pte Ltd, Singapore.}}
\date{2 November 2026\\[3pt]\small
\href{https://creativecommons.org/licenses/by/4.0/}{CC BY 4.0}}
\begin{document}
\maketitle
\begin{abstract}
We resolve the nonlinear classical $k=6$ factor-perfection question of
Pozsgay and Wanless, \S5 immediately after equation (5.13), and their
qubit decomposition conjecture in \S8, for both constructions studied
there ($k=6,8$) \cite{PW}.
For arbitrary maps $A,B,C:X^2\to X^2$ on a finite alphabet, if the graph of
$C^{(23)}B^{(13)}A^{(12)}$ is an index-one orthogonal array of strength three,
then every factor graph has strength two. No invertibility or linearity of
the factors is assumed. A proof using nonnegative integer matrices gives
the stronger equivalence between perfect factors and six specified
permutation flattenings of the composite tensor. A norm--determinant
argument gives the quantum analogue for unitary gates. For arbitrary
complex factors, maximal exterior entanglement instead guarantees a
decomposition by perfect replacement tensors of the same bond dimension.
Consequently no six-qubit absolutely maximally entangled state has this
three-vertex triangular decomposition with bond dimension two. The
eight-qubit case follows from the known nonexistence of an eight-qubit
absolutely maximally entangled state \cite{Rains,Scott}, completing the
qubit conjecture for these two constructions.
\end{abstract}

\section{The question and its classical formulation}
Pozsgay and Wanless \cite[\S5, paragraph after (5.13)]{PW} ask whether
their linear six-leg factor-perfection result extends to nonlinear maps.
We prove the classical assertion for every finite alphabet and give a
quantum replacement theorem that settles their qubit decomposition
conjecture in \S8 for $k=6,8$. Their
factorization is \cite[(4.2)]{PW}
\begin{equation}\label{circuit}
 U=C^{(23)}B^{(13)}A^{(12)}.
\end{equation}
Here the superscripts indicate the two wires on which a map acts.
With exterior symbols $x,y,z,p,q,r$ and interior symbols $u,v,w$, write
\begin{equation}\label{wires}
 A(x,y)=(u,v),\qquad B(u,z)=(p,w),\qquad C(v,w)=(q,r).
\end{equation}
Thus $U(x,y,z)=(p,q,r)$.

For $D=|X|\ge1$, an index-one $\OA(D^s,n,D,s)$ is an array of $D^s$
rows and $n$ columns whose projection onto any $s$ distinct columns is a
bijection onto $X^s$. This is the convention in \cite[\S4]{PW} and
\cite[\S III]{GZ}. The graph of a map $F:X^2\to X^2$ has rows
$(a,b,F_1(a,b),F_2(a,b))$. We call $F$ \emph{perfect} when this graph is
an $\OA(D^2,4,D,2)$; the graph of $U$ is perfect when it is an
$\OA(D^3,6,D,3)$.

These arrays give minimal-support absolutely maximally entangled states
\cite{GZ,GALRZ}. Indeed, if an array has $D^s$ rows on $2s$ columns and
every $s$-column projection is bijective, the normalized sum of its basis
rows has coefficient matrix $D^{-s/2}P$ across every balanced partition,
where $P$ is a permutation matrix. Each reduction is therefore $I/D^s$.
Directional unitarity, studied for other geometries in \cite{MPW}, imposes
a selected family of partitions rather than all balanced partitions.

\begin{theorem}\label{main}
Let $X$ be finite and let $A,B,C:X^2\to X^2$ be arbitrary maps. If the
graph of \eqref{circuit} has strength three and index one, then the graphs
of $A,B,C$ have strength two and index one.
\end{theorem}

\section{A matrix proof and a stronger criterion}
All matrices in the proof are finite. A \emph{flattening} groups a
tensor's indices into an ordered row tuple and an ordered column tuple.
Changing the order within a tuple just permutes matrix rows or columns.

\begin{lemma}\label{matrix}
If square matrices $M_1,\ldots,M_m$ of the same size have nonnegative
integer entries and $M_m\cdots M_1$ is a permutation matrix, then every
$M_j$ is a permutation matrix.
\end{lemma}
\begin{proof}
The product is nonsingular, so every factor is nonsingular. In particular,
no factor has a zero row. For the all-ones column vector $\ind$, integrality
implies $M_j\ind\ge\ind$ entrywise. Nonnegative matrices preserve entrywise
inequalities, and hence
\[
 \ind=M_m\cdots M_1\ind\ge M_m\ind\ge\ind.
\]
Thus every row sum of $M_m$ is one. Each row consists of a single one and
zeros, and nonsingularity makes those ones occupy distinct columns.
Consequently $M_m$ is a permutation matrix. Multiply by $M_m^{-1}$ and
repeat for the remaining product. The case $m=1$ starts the induction.
\end{proof}

Define binary tensors
\[
 a_{xyuv}=\ind_{A(x,y)=(u,v)},\quad
 b_{uzpw}=\ind_{B(u,z)=(p,w)},\quad
 c_{vwqr}=\ind_{C(v,w)=(q,r)}.
\]
Their contraction is
\begin{equation}\label{contract}
 T_{xyzpqr}=\sum_{u,v,w\in X}a_{xyuv}b_{uzpw}c_{vwqr}
             =\ind_{U(x,y,z)=(p,q,r)}.
\end{equation}
The equality holds because an input triple determines exactly one
consistent choice of the interior symbols. It does not presuppose
bijectivity of any gate. A balanced flattening of $T$ is a permutation
matrix exactly when the projections of its support onto both triples
are bijections. In particular, perfection of the composite implies that
every such flattening is a permutation matrix.

For example, $A:xu\to yv$ denotes the matrix with columns indexed by
$(x,u)$, rows indexed by $(y,v)$ and entry $a_{xyuv}$.
An arrow denotes a matrix, not an assumption that the relation is
single-valued in that direction.

\begin{theorem}\label{six}
For $D\ge1$, the following conditions are equivalent:
\begin{enumerate}
\item The three gates $A,B,C$ are perfect.
\item The six flattenings of $T$ whose column triples are
\[
 xyz,\quad xyp,\quad xzp,\quad yzp,\quad xyq,\quad xyr
\]
and whose row triples are the respective complementary triples are
permutation matrices.
\end{enumerate}
Equivalently, the twelve three-column projections containing one of the
pairs $xy$, $zp$, $qr$ are bijective.
\end{theorem}
\begin{proof}
The six flattenings factor as shown in Table~\ref{cuts}. Each factor
acts on two live symbols while the remaining symbol passes unchanged.
It is therefore a row/column permutation of a local $D^2\times D^2$
flattening tensored with $I_D$.

\begin{table}[ht]
\centering\small
\begin{tabular}{@{}ccccc@{}}
\toprule
Exterior input & First factor & Second factor & Third factor & Exterior output\\
\midrule
$xyz$ & $A:xy\to uv$ & $B:uz\to pw$ & $C:vw\to qr$ & $pqr$\\
$xyp$ & $A:xy\to uv$ & $B:up\to zw$ & $C:vw\to qr$ & $zqr$\\
$xzp$ & $B:zp\to uw$ & $A:xu\to yv$ & $C:vw\to qr$ & $yqr$\\
$yzp$ & $B:zp\to uw$ & $A:yu\to xv$ & $C:vw\to qr$ & $xqr$\\
$xyq$ & $A:xy\to uv$ & $C:vq\to wr$ & $B:uw\to zp$ & $zpr$\\
$xyr$ & $A:xy\to uv$ & $C:vr\to wq$ & $B:uw\to zp$ & $zpq$\\
\bottomrule
\end{tabular}
\caption{Six products of lifted local matrices. Factors act left to right
in the table; the matrix product has the reverse written order.}\label{cuts}
\end{table}

To check these products without any invertibility assumption, the live
triples in their six rows are, respectively,
\begin{align*}
xyz&\longrightarrow uvz\longrightarrow pvw\longrightarrow pqr,\\
xyp&\longrightarrow uvp\longrightarrow vzw\longrightarrow zqr,\\
xzp&\longrightarrow xuw\longrightarrow yvw\longrightarrow yqr,\\
yzp&\longrightarrow yuw\longrightarrow xvw\longrightarrow xqr,\\
xyq&\longrightarrow uvq\longrightarrow uwr\longrightarrow zpr,\\
xyr&\longrightarrow uvr\longrightarrow uwq\longrightarrow zpq.
\end{align*}
Each matrix entry is exactly the sum in \eqref{contract}: every interior
symbol is summed once and every exterior symbol remains free.

Assume condition 2. Lemma~\ref{matrix} makes every lifted factor a
permutation matrix. For $D\ge1$, $M\otimes I_D$ is a permutation matrix
if and only if $M$ is: its row and column sums are those of $M$, repeated
$D$ times. Table~\ref{cuts} consequently proves all three balanced local
partitions for each gate:
\[
\begin{array}{c|ccc}
A&xy\mid uv&xu\mid yv&xv\mid yu\\
B&uz\mid pw&up\mid zw&uw\mid zp\\
C&vw\mid qr&vq\mid wr&vr\mid wq.
\end{array}
\]
Transposing a permutation matrix exchanges the two sides. Thus the table
covers all six two-column projections of every gate graph, proving
condition 1.

Conversely, condition 1 makes every local flattening in Table~\ref{cuts}
a permutation matrix. Its lifts and each of the six products are
permutation matrices, proving condition 2. The twelve triples in the
statement are the six column triples and their six distinct complements;
they are exactly the triples containing a full pair among $xy,zp,qr$.
\end{proof}

Theorem~\ref{main} follows since a strength-three graph supplies all
twenty projections and hence the twelve used in Theorem~\ref{six}.
For $X=\varnothing$, every relevant graph and target tuple set is empty;
the required projections are bijective. For $D=1$, all the matrices
above have size one and the argument applies unchanged.

\section{Quantum factors and replacement tensors}
For this section $D\ge1$ and $a,b,c$ in \eqref{contract} may have arbitrary
complex entries. A four-leg tensor is \emph{perfect} if its three
$D^2\times D^2$ balanced flattenings are unitary. A six-leg tensor is
perfect if its ten balanced flattenings are unitary. The corresponding
normalized states have coefficients $D^{-1}a$ and $D^{-3/2}T$,
respectively. Transpose, without complex conjugation, exchanges the two
sides of a flattening and preserves unitarity.

\begin{lemma}\label{normdet}
Let $M_1,\ldots,M_m$ be complex $n\times n$ matrices with
$\|M_j\|_F^2=n$. If $M_m\cdots M_1$ is unitary, every $M_j$ is unitary.
\end{lemma}
\begin{proof}
Writing $s_1,\ldots,s_n$ for the singular values of a factor gives
\[
 |\det M_j|^2=\prod_{i=1}^n s_i^2
 \le \left(\frac{\sum_i s_i^2}{n}\right)^n=1.
\]
The product of these determinant moduli is one, so equality holds for
every factor. Equality in the arithmetic--geometric mean inequality
forces all its squared singular values to be one.
\end{proof}

\begin{theorem}\label{quantum}
If $A,B,C$ are ordinary unitary two-site gates, then their three
four-leg tensors are perfect if and only if the six exterior flattenings
in Table~\ref{cuts} are unitary.
\end{theorem}
\begin{proof}
Every reshuffling of an ordinary unitary gate has squared Frobenius norm
$D^2$. Its spectator lift has size $D^3$ and squared Frobenius norm
$D^3$. The products in Table~\ref{cuts} hold for complex entries as well.
Lemma~\ref{normdet} therefore makes every local balanced flattening
unitary. Conversely, perfect factors give products of unitaries in all
six rows, as established by Pozsgay and Wanless
\cite[\S3.2, Theorems 1--2]{PW}.
\end{proof}

For arbitrary complex factors the supplied tensors need not be perfect:
an invertible matrix on one end of an internal bond and its inverse on
the other end leave the contraction unchanged. The correct conclusion
is the existence of perfect replacement tensors. The next theorem gives
a converse at the level of decompositions to the forward construction
of \cite[\S3.2, Theorems 1--2]{PW}.

\begin{theorem}\label{replacement}
Suppose a complex six-leg tensor has the triangular contraction
\eqref{contract}, with every exterior and interior index of dimension
$D$. If its six flattenings in Table~\ref{cuts} are unitary, the same
tensor has a contraction of this form with three perfect four-leg
tensors and the same bond dimension $D$.
\end{theorem}
\begin{proof}
Group the exterior legs as $a=(x,y)$, $b=(z,p)$ and $c=(q,r)$, and put
\[
 |\Omega\rangle=D^{-3/2}\sum_{u,v,w}
 |u,v\rangle_a|u,w\rangle_b|v,w\rangle_c.
\]
This is a product of three normalized maximally entangled pairs; each
grouped marginal is $I_{D^2}/D^2$. Define square local matrices by
\[
 (F_a)_{xy,uv}=a_{xyuv},\quad
 (F_b)_{zp,uw}=b_{uzpw},\quad
 (F_c)_{qr,vw}=c_{vwqr}.
\]
Then the normalized exterior state is
$|\Psi\rangle=(F_a\otimes F_b\otimes F_c)|\Omega\rangle$.
The unitary flattenings imply that $\Psi$ is normalized and each grouped
marginal is $I_{D^2}/D^2$: the pairs $xy$ and $qr$ lie on sides of
$xyz\mid pqr$, and $zp$ lies on a side of $xzp\mid yqr$.
Full rank of these marginals makes all three $F_i$ invertible.

Write the polar decompositions $F_i=V_iP_i$, with $V_i$ unitary and
$P_i$ positive definite. Let
\[
 H=(\log P_a)\otimes I\otimes I+
 I\otimes(\log P_b)\otimes I+
 I\otimes I\otimes(\log P_c).
\]
The state $|\Phi\rangle=(V_a^\dagger\otimes V_b^\dagger\otimes
V_c^\dagger)|\Psi\rangle=e^H|\Omega\rangle$ is also normalized with
the same grouped marginals. For $f(t)=\|e^{tH}\Omega\|^2$, these
marginals give
\[
 f'(0)=\frac{2}{D^2}\sum_{i=a,b,c}\operatorname{tr}(\log P_i)
       =f'(1),\qquad
 f''(t)=4\|He^{tH}\Omega\|^2\ge0.
\]
The integral of $f''$ on $[0,1]$ vanishes. By continuity $H\Omega=0$,
so $e^H\Omega=\Omega$ and
\[
 |\Psi\rangle=(V_a\otimes V_b\otimes V_c)|\Omega\rangle.
\]
This is precisely the same triangular tensor, with the entries of the
three $V_i$ as replacement vertices. Each has squared Frobenius norm
$D^2$ in every reshuffling. Lemma~\ref{normdet} applied to the six
unchanged exterior flattenings now makes every balanced local
flattening unitary.
\end{proof}

The polar argument is the equal-marginal critical-state uniqueness
principle underlying Kempf--Ness theory; compare
\cite[Theorem 2 and the following paragraph]{GKW}. The proof above
establishes the needed statement directly for grouped dimension $D^2$.
It also shows that any triangle state of this form with maximally mixed
grouped marginals on $xy$, $zp$ and $qr$ is local-unitarily equivalent,
with respect to these three dimension-$D^2$ nodes, to three maximally
entangled pairs of dimension $D$; at $D=2$ these are three Bell pairs.

\begin{corollary}\label{qubits}
No six-qubit absolutely maximally entangled state admits the triangular
three-vertex contraction \eqref{contract} with interior bond dimension
two, even with arbitrary complex starting tensors.
\end{corollary}
\begin{proof}
Such a state would satisfy all ten balanced unitary conditions after
the stated normalization. Theorem~\ref{replacement} would then produce
perfect four-qubit tensors, which do not exist
\cite[\S2]{HS}.
\end{proof}

Pozsgay and Wanless \cite[\S8]{PW} write that they ``conjectured that
AME states with qubits cannot be decomposed in the way that we
formulated here''; see also the end of \S3.2. This conjecture is now
settled completely for both constructions in their paper: Corollary~\ref{qubits}
handles $k=6$, with the stated topology and bond dimension, while $k=8$
is vacuous because no eight-qubit AME state exists
\cite[Theorem 14 in the arXiv version]{Rains}\cite[\S III, after Table I]{Scott}.
Indeed, such an AME state would define a pure $((8,1,5))$ quantum code,
whereas Rains's bound for length $6m+\ell=8$ ($m=1$, $\ell=2$) gives
$d\le 2m+2=4$.
Kraft et al.\ \cite[Table I]{Kraft}
previously studied triangle-network entanglement and obtained a
putative squared maximal overlap $1/2$ for a six-qubit AME target by a
see-saw algorithm. Their optimization concerns pair sources and local
unitaries. Here the exclusion is analytic and allows arbitrary complex
vertices with the specified bond dimension; no overlap optimum is asserted.

\section{Consequences and examples}
\begin{corollary}\label{reverse}
Theorem~\ref{main} holds for every chronological order of the three
distinct pair interactions $12,13,23$, including the opposite order
$U=A^{(12)}B^{(13)}C^{(23)}$.
\end{corollary}
\begin{proof}
Permuting three wires induces every permutation of the three pair
interactions. Apply the original theorem after such a relabeling;
swapping the arguments and outputs of an individual gate preserves
perfection. In particular, under $1\leftrightarrow3$ the gates act, in
chronological order, on $(1,2),(1,3),(2,3)$. The relabeling may exchange
the two arguments and the two outputs of a gate; such exchanges simply
permute columns of its graph. The same is true of the six exterior
columns. Apply Theorem~\ref{main} and undo those permutations.
\end{proof}

The forward implication of Theorem~\ref{six}, from perfect factors to
twelve orthogonal projections, is the classical observation of
\cite[\S4, following (4.2)]{PW}. The converse supplies an exact criterion.
The other eight triples choose one symbol from each of $xy,zp,qr$.
Their orthogonality is an additional condition, as the following example
shows.

Over $\mathbb F_3$ take all three gates equal to
\[
 F(a,b)=(a+b,a-b).
\]
Every coefficient is nonzero and the determinant is $-2\ne0$, so the
two outputs are Latin in each argument and $F$ is bijective; these are
precisely the six two-column conditions. Yet the composite is
\[
 U(x,y,z)=(x+y+z,\,2x-z,\,-2y+z).
\]
For fixed $x,z$, the second output is independent of $y$. The projection
onto $(x,z,q)$ is not injective, so the composite is not perfect. This is
the four-column array of \cite[Table 1]{PW} used at each gate.

\begin{corollary}\label{dual}
For classical gates, the three flattenings
$xyz\mid pqr$, $yzp\mid xqr$ and $xyr\mid zpq$ are permutation
matrices if and only if the following two flattenings of each gate are
permutation matrices:
\[
\begin{array}{c|cc}
A&xy\mid uv&xv\mid yu\\
B&uz\mid pw&uw\mid zp\\
C&vw\mid qr&vr\mid wq.
\end{array}
\]
For ordinary unitary quantum gates, the same equivalence holds with
``unitary'' in place of ``permutation''.
\end{corollary}
\begin{proof}
Rows 1, 4 and 6 of Table~\ref{cuts} use precisely these local
partitions, with some sides exchanged. Apply Lemma~\ref{matrix} in the
classical case and Lemma~\ref{normdet} in the quantum case. The converse
in each case follows by multiplication.
\end{proof}
These are the ordinary and geometrically dual local partitions and the
three hexagonal exterior partitions in \cite[\S3.2, Theorem 1]{PW}.
Thus the corollary supplies the converse for the specified triangular
factorization.

\paragraph{Irredundancy of the six specified conditions.}
None of the six conditions in Theorem~\ref{six} can be deleted from
that family. To see this, over $\mathbb F_3$ use $F$ at all gates
except the replacement in the following table. The last column gives
the numbers of distinct projections on the two sides of the indicated
partition. Every other named partition has $27$ projections on each
side, so its flattening is a permutation matrix.
\[
\begin{array}{c|c|c}
\text{failed partition}&\text{replacement}&\text{projection sizes}\\\hline
xyz\mid pqr&B(u,z)=(u+z,u+z)&27,9\\
xyp\mid zqr&B(u,z)=(u,z)&9,9\\
xzp\mid yqr&A(x,y)=(x,y)&9,9\\
yzp\mid xqr&A(x,y)=(y,x)&9,9\\
xyq\mid zpr&C(v,w)=(v,w)&9,9\\
xyr\mid zpq&C(v,w)=(w,v)&9,9
\end{array}
\]
The sizes follow by taking ranks of the two $3\times3$ linear
projection maps over $\mathbb F_3$; they are also enumerated in
\texttt{verify.py}. This establishes irredundancy of the named family,
not minimality among all possible characterizations.

\paragraph{Nonlinear factors.}
Internal symbol permutations preserve both the composite and local
perfection. If $\phi,\psi,\theta$ are permutations of $X$, replace the
gates by
\begin{align*}
 A'(x,y)&=(\phi(A_1(x,y)),\psi(A_2(x,y))),\\
 B'(u,z)&=(B_1(\phi^{-1}(u),z),
                 \theta(B_2(\phi^{-1}(u),z))),\\
 C'(v,w)&=C(\psi^{-1}(v),\theta^{-1}(w)).
\end{align*}
Substitution in \eqref{wires} proves equality of the composites.
Each local graph changes by separate symbol permutations in its columns,
so its orthogonal-array property is unchanged.
For example, use the three matrices of \cite[Example 3]{PW} over
$\mathbb F_5$,
\[
 A=\begin{pmatrix}3&4\\3&1\end{pmatrix},\quad
 B=\begin{pmatrix}4&1\\1&3\end{pmatrix},\quad
 C=\begin{pmatrix}3&1\\3&4\end{pmatrix},\quad
 G=\begin{pmatrix}2&1&1\\2&2&3\\1&4&2\end{pmatrix}.
\]
The choice $\phi(t)=\psi(t)=\theta(t)=t^3$, a self-inverse permutation
of $\mathbb F_5$, gives three non-affine gates with the same perfect
composite $G$. This example is a relabeling of the cited linear example,
not a new equivalence class of arrays.

To preserve the Yang--Baxter equality of the two orders as well, conjugate
every wire by the same symbol permutation $s$. Set
\[
 A_s=(s\times s)A(s^{-1}\times s^{-1}),
 \quad B_s=(s\times s)B(s^{-1}\times s^{-1}),
 \quad C_s=(s\times s)C(s^{-1}\times s^{-1}).
\]
Writing $S=s\times s\times s$, either ordered product is the
corresponding original product conjugated by $S$. Example 3 of
\cite{PW} has equal products in the two orders; hence the equality
survives. For $s(t)=t^3$ the three conjugated gates are non-affine,
perfect, and have a perfect composite. The exact finite checks enumerate
all their rows. These are non-affine representatives of the same
symbol-equivalence class.

\section{Coverage and mathematical verification}
\begin{table}[ht]
\centering\small
\begin{tabular}{@{}p{.35\linewidth}p{.60\linewidth}@{}}
\toprule
Item in Pozsgay--Wanless \cite{PW} & Coverage\\
\midrule
\S5, after (5.13): nonlinear classical $k=6$ factors
 & Proved by Theorem~\ref{main}.\\
\S3.1: opposite triangular order
 & All six chronological orders covered by Corollary~\ref{reverse}.\\
\S4: twelve orthogonal projections from perfect factors
 & Prior implication reproduced and strengthened to an equivalence in
   Theorem~\ref{six}.\\
\S3.2: perfection of all factors versus all ten balanced partitions
 & Converse fails; the $\mathbb F_3$ example separates the conditions.\\
\S3.2, Theorems 1--2: dual/perfect factors give three/six unitary flattenings
 & Prior forward implications credited; converses proved in
   Corollary~\ref{dual} and Theorem~\ref{quantum}.\\
\S3.1 and Example 3: perfect Yang--Baxter factors
 & Prior existence example cited; global wire conjugation gives
   non-affine representatives preserving both orders.\\
\S\S5--6: linear existence tests and small alphabet cases
 & Existing linear conclusions retained; nonlinear existence at $D=4$
   is not determined by the factor-perfection implication.\\
\S7: eight-leg networks at $D\ge3$
 & Not covered. Example 5 already has an identity factor and a perfect
   composite, so an unrestricted eight-leg analogue is false.\\
\S8 (also end of \S3.2): qubit decomposition conjecture, $k=6,8$
 & Settled completely: $k=6$ by Corollary~\ref{qubits}, including arbitrary
   complex vertices of bond dimension two; $k=8$ by nonexistence of
   eight-qubit AME states \cite{Rains,Scott}.\\
\S8: quantum $D=3,4$ existence; Appendix B: graph-state gate counts
 & Existence is reduced to perfect replacements; these existence
   questions and unrestricted gate counts remain undetermined here.\\
\bottomrule
\end{tabular}
\caption{Coverage of the posing question and neighboring statements.}
\label{coverage}
\end{table}

\paragraph{Mathematical verification.}
The declaration \texttt{result : claim} in the pinned
\href{https://github.com/the-omega-institute/trureturing/blob/6715a5bea38a03c80f6c4da6de93fea88798f026/D5/S3/Combinatorics/Latin/ClassicalPerfectTensorFactorization.lean}{Lean source}
belongs to the namespace
\begin{center}\small
\nolinkurl{D5.S3.Combinatorics.Latin.ClassicalPerfectTensorFactorization}.
\end{center}
It proves precisely
Theorem~\ref{main}'s implication for every finite alphabet, including the
empty case, with \texttt{circuitGraph} columns $(x,y,z,p,q,r)$ and
\texttt{gateGraph} columns consisting of the two inputs and two outputs.
All other results, including the quantum replacement theorem and its
six-qubit consequence, have written proofs here; the eight-qubit
nonexistence statement is the cited consequence of Rains's bound.
The accompanying source
archive contains \texttt{verify.py} and its deterministic expected output
\texttt{verification.json}. It checks the displayed examples by exact
arithmetic, exhausts all $24^3$ binary permutation-gate circuits, and
checks Corollary~\ref{dual} on all $81^3$ triples of linear gates over
$\mathbb F_3$. These finite checks supplement the general proofs.

\begin{thebibliography}{9}
\bibitem{PW}
B. Pozsgay and I. M. Wanless,
\emph{Tensor network decompositions for absolutely maximally entangled
states}, Quantum \textbf{8} (2024), 1339.
\href{https://doi.org/10.22331/q-2024-05-08-1339}{doi:10.22331/q-2024-05-08-1339}.
\href{https://arxiv.org/abs/2308.07042v3}{arXiv:2308.07042v3}.

\bibitem{GZ}
D. Goyeneche and K. \.{Z}yczkowski,
\emph{Genuinely multipartite entangled states and orthogonal arrays},
Phys. Rev. A \textbf{90} (2014), 022316.
\href{https://doi.org/10.1103/PhysRevA.90.022316}{doi:10.1103/PhysRevA.90.022316}.

\bibitem{GALRZ}
D. Goyeneche, D. Alsina, J. I. Latorre, A. Riera, and K. \.{Z}yczkowski,
\emph{Absolutely maximally entangled states, combinatorial designs, and
multiunitary matrices}, Phys. Rev. A \textbf{92} (2015), 032316.
\href{https://doi.org/10.1103/PhysRevA.92.032316}{doi:10.1103/PhysRevA.92.032316}.

\bibitem{MPW}
M. Mesty\'an, B. Pozsgay, and I. M. Wanless,
\emph{Multi-directional unitarity and maximal entanglement in spatially
symmetric quantum states}, SciPost Phys. \textbf{16} (2024), 010.
\href{https://doi.org/10.21468/SciPostPhys.16.1.010}{doi:10.21468/SciPostPhys.16.1.010}.

\bibitem{GKW}
G. Gour, B. Kraus, and N. R. Wallach,
\emph{Almost all multipartite qubit quantum states have trivial stabilizer},
J. Math. Phys. \textbf{58} (2017), 092204.
\href{https://doi.org/10.1063/1.5003015}{doi:10.1063/1.5003015}.

\bibitem{HS}
A. Higuchi and A. Sudbery,
\emph{How entangled can two couples get?}, Phys. Lett. A
\textbf{273} (2000), 213--217.
\href{https://doi.org/10.1016/S0375-9601(00)00480-1}{doi:10.1016/S0375-9601(00)00480-1}.

\bibitem{Kraft}
T. Kraft, S. Designolle, C. Ritz, N. Brunner, O. G\"uhne, and M. Huber,
\emph{Quantum entanglement in the triangle network}, Phys. Rev. A
\textbf{103} (2021), L060401.
\href{https://doi.org/10.1103/PhysRevA.103.L060401}{doi:10.1103/PhysRevA.103.L060401}.

\bibitem{Rains}
E. M. Rains,
\emph{Quantum shadow enumerators}, IEEE Trans. Inf. Theory
\textbf{45} (1999), 2361--2366.
\href{https://doi.org/10.1109/18.796376}{doi:10.1109/18.796376}.
\href{https://arxiv.org/abs/quant-ph/9611001v2}{arXiv:quant-ph/9611001v2}.

\bibitem{Scott}
A. J. Scott,
\emph{Multipartite entanglement, quantum-error-correcting codes, and
entangling power of quantum evolutions}, Phys. Rev. A
\textbf{69} (2004), 052330.
\href{https://doi.org/10.1103/PhysRevA.69.052330}{doi:10.1103/PhysRevA.69.052330}.
\href{https://arxiv.org/abs/quant-ph/0310137}{arXiv:quant-ph/0310137}.
\end{thebibliography}
\end{document}
